\section{Introduction and objectives}
\label{sec:intro}
ThermalSense is a real-time 3D scene of a forward military outpost. It started as a lab exercise: build objects only
from unit primitives with a scale, rotation and translation table. It grew into a full renderer and a small game.
This report follows the graphics pipeline in the order taught in class \cite{hearn2011,foley1996} (Part~I), then shows the advanced work
built on it (Part~II) and closes with the measured results (Part~III).

\textbf{Objectives.}
\begin{enumerate}[label=O\arabic*.]
\item Model a complete scene from a small set of primitives with hierarchical transforms (Sections~\ref{sec:prims}--\ref{sec:hier}, \ref{sec:decomp}).
\item Implement and visualise each pipeline stage: viewing, projection, clipping, hidden-surface removal, rasterisation, illumination, shading, texturing (Sections~\ref{sec:view}--\ref{sec:tex}).
\item Show ray-tracing concepts inside the same scene (Sections~\ref{sec:rt}, \ref{sec:rtadv}).
\item Drive lighting from a physical day/night model (Section~\ref{sec:day}).
\item Build a simulated thermal sensor and an image-processing chain on its output (Section~\ref{sec:thermal}).
\item Make the scene interactive: controllable units, \ab{471}{AI}, path planning (Sections~\ref{sec:dynamic}, \ref{sec:game}, \ref{sec:paths}).
\item Measure performance and state the trade-offs (Sections~\ref{sec:hard}, \ref{sec:perf}, \ref{sec:limits}).
\end{enumerate}

\textbf{What is simulated.} The ``thermal'' image is not real sensing. Each part carries a temperature value; a
sensor model blurs, adds noise and fixed-pattern error, and maps heat to a false-colour palette
(Section~\ref{sec:thermal}). The report describes only features that work, fully or in part.
